% !TeX program = xelatex
\documentclass{cjlu_opening}

\usepackage{array}
\usepackage{longtable}
\usepackage{tabularx}

\studentname{李康峰}
\studentid{2201400216}
\major{自动化}
\classinfo{22 工试\hspace{0.5em}2 班}
\thesistitle{基于人工智能的缺陷射线图像识别与检测系统设计}
\advisor{洪宇翔}
\school{机电工程学院}
\coverdate{2025 年 12 月 30 日}
\covermaintitle{文献综述}

\begin{document}

\makeopeningcover

文献综述能较全面地反映与本课题直接相关的国内外研究成果，特别是近年来的最新成果和发展趋势，
能指出该课题所需要进一步解决的问题。

文献综述内容需体现：1、综合性。要尽可能把所有重要研究成果搜集到手，使各种派别的观点清楚明晰，
不遗漏重要的派别和观点。2、描述性。要求对被介绍的观点作客观性的描述。在归纳各种观点时要抓住要点，
表述时应简明扼要。3、评价性。要对各种成果进行恰当而中肯的评价，并表明作者自己的观点和主张。
文献综述重点在于“述”，要点在于“评”。

文献综述主体部分的结构，一般应包括该课题的“研究历史”的回顾、“研究现状”的对比，以及研究的“发展趋势”。
文献综述的格式规范与正文的通用格式一致。一般地说，为撰写本科毕业设计（论文）提供观点和材料基础的文献综述字数至少要在 3000 字以上。

\par\noindent{\heiti\zihao{-4}参考文献}\par
要求不少于 20 篇期刊文献（英文文献不少于 5 篇，且必须包含要翻译的 2 篇英文文献）。

\renewcommand{\arraystretch}{1.25}
\setlength{\tabcolsep}{6pt}
\begin{longtable}{|>{\centering\arraybackslash}m{0.18\textwidth}|m{0.74\textwidth}|}
    \hline
    {\heiti 文献类型}      & {\heiti 格式与示例}                                                                                                  \\
    \hline
    专著、论文集、学位论文、报告和会议录 &
    格式：[序号]主要责任者．文献题名[文献类型标识]．版次（第 1 版可省略）．出版地:出版者,出版年．起止页码。                                                                             \\
                       & 示例：[1]刘国钧等．图书馆目录[M]．北京:高等教育出版社,1957．15-18。                                                                      \\
    \hline
    期刊文章               &
    格式：[序号]主要责任者．文献题名[文献类型标识]．刊名,年,卷(期):起止页码。                                                                                            \\
                       & 示例：[2]金显贺,王昌长,王忠东等．一种用于在线检测局部放电的数字滤波技术[J]．清华大学学报(自然科学版),1993,33(4):62-67。                                       \\
    \hline
    论文集中的析出文献          &
    格式：[序号]析出文献主要责任者．析出文献题名[文献类型标识]．原文献主要责任者．原文献题名[C]．出版地:出版者,出版年．析出文献起止页码。                                                              \\
                       & 示例：[3]钟文发．非线性规划在可燃毒物配置中的应用[A]．赵玮．运筹学的理论与应用——中国运筹学会第五届大会论文集[C]．西安:西安电子科技大学出版社,1996．468-471。                      \\
    \hline
    国际、国家标准名称          &
    格式：[序号]标准编号,标准名称[文献类型标识]。                                                                                                            \\
                       & 示例：[4]GB/T 16159-1996,汉语拼音正词法基本规则[S]。                                                                           \\
    \hline
    专利                 &
    格式：[序号]专利所有者．专利题名[文献类型标识]．专利国别:专利号,出版日期。                                                                                             \\
                       & 示例：[5]姜锡洲．一种温热外敷药制备方案[P]．中国专利:881056073,1989-07-26。                                                             \\
    \hline
    报纸文章               &
    格式：[序号]主要责任者．文献题名[文献类型标识]．报纸名,出版日期（版次）。                                                                                              \\
                       & 示例：[6]谢希德．创造学习的新思路[N]．人民日报,1998-12-25(10)。                                                                      \\
    \hline
    电子文献               &
    格式：[序号]主要责任者．电子文献题名[电子文献及载体类型标识]．电子文献出处或可获得地址,发表或更新日期/引用日期（任选）。                                                                      \\
                       & 示例：[7]王明亮．关于中国学术期刊标准化数据库系统工程的进展[EB/OL]．http://www.cajcd.edu.cn/pub/wml.txt/980810-2.html,1998-08-16/1998-10-04。 \\
    \hline
\end{longtable}

\end{document}
